% первая часть

\section{Теоретические основы и анализ существующих решений}
\subsection{Интеллектуальные обучающие системы: структура и компоненты}
Интеллектуальная обучающая система (ИОС) — это программный комплекс, 
обеспечивающий автоматизированное управление процессом обучения с 
учётом индивидуальных особенностей пользователя~\cite{brusilovsky2001adaptive}. В отличие от 
традиционных систем управления обучением, ИОС обладает 
механизмами адаптации, диагностики знаний и прогнозирования результатов~\cite{brusilovsky2001adaptive}.

Архитектура ИОС обычно включает в себя четыре ключевых компонента.
Первый — модель предметной области, содержащая формализованное 
представление знаний по изучаемой дисциплине: онтологии, семантические 
графы, структурированные базы заданий. Второй — модель обучающегося, 
динамически обновляемое представление текущего состояния знаний, навыков и 
характеристик студента~\cite{brusilovsky2007usermodels}. Третий — педагогическая модель, определяющая стратегии 
выбора учебного контента, формирования последовательности заданий и выдачи 
обратной связи. Четвёртый — интерфейс взаимодействия, обеспечивающий удобный 
доступ студента и преподавателя к функциям системы.
В современных ИОС эта классическая схема расширяется за счёт дополнительных 
компонентов. Модуль аналитики отслеживает поведенческие паттерны студента: 
время выполнения заданий, типичные ошибки, ритм сессий обучения~\cite{brusilovsky2007usermodels}. 
Модуль геймификации преобразует учебные действия в игровые события — 
начисление баллов, разблокировку достижений, позиционирование в рейтинге —
что повышает внутреннюю мотивацию~\cite{dominguez2013gamifying, plass2015foundations}. Рекомендательная подсистема на основе 
накопленных данных предлагает следующие шаги обучения, учитывая как текущий 
уровень знаний, так и контекстные факторы~\cite{brusilovsky2007usermodels, conati2002bayesian}.

Учебный контент в ИОС организован иерархически: курс → модуль → тема → задание. 
Каждый элемент этой иерархии снабжается метаданными: уровень сложности, 
предварительные требования (пресуппозиции), компетенции, формируемые после 
освоения. Это позволяет строить граф зависимостей знаний и автоматически подбирать
оптимальную последовательность изучения материала для конкретного студента~\cite{woolf2009building}.

Оценка знаний в ИОС реализуется на нескольких уровнях:
формативная оценка (промежуточная) обеспечивает непрерывный мониторинг 
прогресса и запускает адаптацию траектории в режиме реального времени,
суммативная оценка (итоговая) фиксирует финальный уровень компетентности,
диагностическая оценка (входная) определяет начальный уровень знаний и 
позволяет обойти уже освоенный материал~\cite{brusilovsky2007usermodels, vanlehn2006behavior}.

\subsection{Модели построения когнитивного профиля}
Адаптивное обучение предполагает динамическое управление учебным процессом 
на основе непрерывно обновляемой информации об обучающемся. 
Ключевым инструментом такого управления является когнитивный профиль — 
формализованная модель психологических и академических характеристик студента~\cite{brusilovsky2007usermodels}.

В разрабатываемой системе когнитивный профиль строится на основе четырёх факторов, 
получаемых в ходе психологического тестирования при регистрации и обновляемых в
контрольных точках. Фактор A1 (самоорганизация) оценивает способность студента 
к планированию учебного времени, структурированию задач и самодисциплине. 
Студенты с высоким A1 успешно работают в асинхронном режиме и не нуждаются во 
внешних напоминаниях; для студентов с низким A1 система генерирует 
структурированные расписания и промежуточные дедлайны. Фактор A2 (мотивация) 
разделяет внутреннюю мотивацию (интерес к предмету) и внешнюю 
(стремление к оценке, признанию). Для студентов с преобладанием внешней
мотивации особую роль играют геймификационные механики — лидерборды, 
значки достижений, прогресс-бары~\cite{dominguez2013gamifying, bartle1996players}. Фактор A3 (стратегия обучения) 
определяет предпочтительный модальный стиль усвоения информации: визуальный 
(диаграммы, схемы, видео), аудиальный (лекции, подкасты), кинестетический 
(практические задания, симуляции). Система адаптирует формат подачи контента 
в соответствии с доминирующей стратегией~\cite{brusilovsky2007usermodels}. Фактор A4 (интегральный показатель) 
является обобщённой оценкой, синтезирующей A1–A3 и отражающей общую готовность 
студента к эффективному самообучению~\cite{brusilovsky2007usermodels}.

На основе комбинации значений этих факторов система формирует три типа 
рекомендаций. Рекомендации для преподавателя (тип T) описывают оптимальный 
формат взаимодействия с данным студентом: частоту проверок, уровень 
детализации обратной связи, степень автономии. Рекомендации по обратной 
связи (тип R) задают параметры автоматических откликов системы: мгновенные 
подсказки или отложенное объяснение, акцент на ошибках или на достижениях. 
Рекомендации для учебного сообщества (тип C) определяют, насколько студенту 
показаны совместные форматы работы: групповые проекты, парное программирование, 
дискуссионные форумы~\cite{brusilovsky2007usermodels, koper2005learner}.

\subsection{Контекстная адаптация в ИОС}
Контекстная адаптация — это корректировка поведения системы на основе информации 
о внешних условиях функционирования, не зависящих непосредственно от учебного 
контента. В разрабатываемой системе реализована контекстная адаптация по трём параметрам: 
географическому положению, погодным условиям и температуре воздуха~\cite{chen2008context, ipapi2026, openweathermap2026}.

Научное обоснование связи погодных условий с когнитивной продуктивностью восходит к исследованиям в 
области экологической психологии. При температуре ниже 13~°C 
или выше 33~°C наблюдается статистически значимое снижение точности 
выполнения когнитивных задач и скорости обработки информации~\cite{keller2005warm}. Дождь и пасмурность 
снижают уровень серотонина, что отрицательно сказывается на концентрации внимания 
и рабочей памяти — механизм, хорошо изученный в контексте сезонных аффективных 
расстройств~\cite{keller2005warm}.

Контекстные данные в системе используются для расчёта итогового параметра 
активности пользователя $P$ по формуле:
\begin{equation}
    \label{eq:P}
    P = K_{\text{region}} \cdot K_{\text{temp}} \cdot K_{\text{weather}}
\end{equation}

где $K_{\text{region}}$ — региональный коэффициент (учитывает климатические особенности), 
$K_{\text{weather}}$ — погодный коэффициент (нормализован к шкале 0–1), $K_{\text{temp}}$ — температурный 
коэффициент (отклонение от комфортного диапазона 15–25~°C). 
Значение $K_{\text{temp}}$ вычисляется как:
\begin{equation}
    \label{eq:Ktemp}
    K_{\text{temp}} = \frac{T_{\text{current}} - T_{\text{comfort}}}{T_{\text{max.diff}}}
\end{equation}
где
\begin{itemize}
    \item $T_{\text{comfort}}$ — комфортная температура (в системе 20~°C);
    \item $T_{\text{current}}$ — текущая температура;
    \item $T_{\text{max.diff}}$ — максимальное допустимое отклонение температуры, при котором 
    коэффициент снижается до минимального значения.
\end{itemize}
    
Выбор диапазона 15–25~°C как комфортного опирается на исследования терморегуляции 
и умственной работоспособности: при температуре ниже 13~°C или выше 33~°C 
наблюдается статистически значимое снижение точности выполнения когнитивных     
задач~\cite{keller2005warm}. Значение $T_{\text{comfort}} = 20$~°C соответствует середине этого диапазона и 
является стандартным нормативом для учебных помещений согласно ГОСТ~30494-2011. 
Нижняя граница $K_{\text{temp}} = 0{,}1$ (а не 0) введена намеренно: полное обнуление 
параметра $P$ привело бы к полному отключению рекомендательной системы, что
нецелесообразно — даже в экстремальных условиях система должна продолжать 
работу в минимальном режиме~\cite{chen2008context}.

Значения $K_{\text{weather}}$ определены на основе анализа психофизиологических исследований 
влияния погодных условий на когнитивную продуктивность. Ясная погода ($K = 1{,}0$) 
принята за базовый уровень, соответствующий оптимальным условиям концентрации 
внимания. Облачность ($K = 0{,}9$) оказывает минимальное прямое воздействие, однако 
незначительно снижает уровень естественной освещённости, что влияет на циркадные 
ритмы. Дождь ($K = 0{,}7$) создаёт акустические помехи и ассоциируется со снижением 
уровня серотонина, что подтверждается исследованиями сезонных аффективных 
расстройств~\cite{keller2005warm}. Гроза ($K = 0{,}5$) формирует выраженный стрессовый фон, 
активирующий симпатическую нервную систему и существенно снижающий способность 
к длительной концентрации. Снег ($K = 0{,}6$) помимо психологического дискомфорта 
создаёт практические препятствия для очного обучения и увеличивает когнитивную 
нагрузку, связанную с логистикой~\cite{chen2008context}.

\subsection{Полный анализ существующей системы}
\subsubsection{Функциональный состав текущей версии}
На момент начала преддипломной практики система реализована в виде 
веб-приложения на Django~\cite{django2026docs} и включает следующие модули:
\begin{itemize}
     \item компонент регистрации и аутентификации пользователей;
     \item база данных с учебными материалами, тестами и результатами прохождения;
     \item рекомендательная система, учитывающая погодные условия при подборе контента~\cite{chen2008context};
     \item интерфейс для преподавателей и администраторов;
     \item модуль формирования личностного портрета обучающегося (в стадии разработки).
\end{itemize}

\subsubsection{Оценка интеллектуальных возможностей системы}
В текущей версии системы реализованы базовые функции, 
характерные для современных ИОС~\cite{woolf2009building, brusilovsky2001adaptive}:
\begin{itemize}
     \item прохождение учебных курсов, структурированных по темам и заданиям;
     \item система оценки с выдачей результатов в виде баллов и общих рекомендаций;
     \item интерфейс для преподавателя для отслеживания прогресса студентов;
     \item персонализация — учёт индивидуальных целей и темпа построен 
     исключительно на результатах тестов и базовой погодной рекомендации, 
     без учёта всего когнитивного профиля~\cite{brusilovsky2007usermodels};
\end{itemize}

\subsection{Вывод по анализу}
В нынешнем виде система располагает устойчивой архитектурной базой, пригодной для 
масштабирования и внедрения более продвинутых образовательных моделей~\cite{gamma1994design}.
Интеллектуальные функции реализованы на базовом уровне и нуждаются в развитии. 
Имеются все предпосылки для внедрения адаптивных механизмов, алгоритмов машинного 
обучения и интеллектуальной диагностики знаний~\cite{ifenthaler2020learning, vanlehn2006behavior}. Полученная картина состояния системы
позволяет перейти к формулированию требований к новой модели обучения.